%------------------------
% Resume Template
% Author : Harsh
% Track  : AI4Engineering Science (Mistral AI, "AI Scientist - Physics Models")
%------------------------
%
% ==== OPEN ITEMS — fill these in before sending, do not guess ====
% 1. [DONE] Expected PhD submission/defence date -- filled: 2026 (expected).
% 2. Leakage-onset framework: dataset scale — number of microstructure realizations / load
%    cases simulated, and turnaround vs. full-scale FEM if you have it
%    (search "[x] full-scale FEM runs")
% 3. Monte Carlo UQ: number of samples run and/or the size of the probability-of-failure
%    study (search "[x] Monte Carlo samples") -- not yet added as a placeholder below;
%    add a number here if you have one.
% 4. [DONE] Physics-based CNN / agentic app: corpus size and validation accuracy filled
%    from the same real figures already used in CV_EN_BMW_AgenticAI.tex (50 materials
%    documents, 93% faithfulness, 100+ held-out cases).
% 5. [DONE] CT-scan defect-detection metrics filled from the same real figures already
%    used in CV_EN_BMW_AgenticAI.tex (91% accuracy, 89 precision / 91 recall, 1200 CT
%    scans, <0.5 sec/image review time).
% 6. If you have used any specific ML architecture keywords relevant to this role (e.g.
%    graph neural networks, neural operators, transformers on mesh/point-cloud data,
%    multi-fidelity training) beyond what's listed here, add them — Mistral's posting
%    explicitly asks about training-strategy and scaling-behavior experience.
% ===================================================================

\documentclass[a4paper,20pt]{article}

\usepackage{latexsym}
\usepackage[empty]{fullpage}
\usepackage{titlesec}
\usepackage{marvosym}
\usepackage[usenames,dvipsnames]{color}
\usepackage{verbatim}
\usepackage{enumitem}
\usepackage[pdftex]{hyperref}
\usepackage{fancyhdr}
\usepackage{graphicx}
\usepackage{tikz}
\usepackage{blindtext}
\usepackage[document]{ragged2e}
\usepackage{xcolor}
\usepackage{hhline}
\usepackage{hyperref}
\usepackage{xcolor}

\hypersetup{
    colorlinks=true,
    linkcolor=blue,
    urlcolor=blue,
    citecolor=blue
}
\usepackage{tikz}

\newcommand{\skilllevel}[1]{
  \begin{tikzpicture}[baseline=-0.5ex]
    \foreach \i in {1,...,5} {
      \ifnum\i>#1
        \draw[fill=white] (\i*0.25,0) circle (0.08);
      \else
        \draw[fill=black] (\i*0.25,0) circle (0.08);
      \fi
    }
  \end{tikzpicture}
}

\pagestyle{fancy}
\fancyhf{}
\fancyfoot{}
\renewcommand{\headrulewidth}{0pt}
\renewcommand{\footrulewidth}{0pt}

\addtolength{\oddsidemargin}{-0.530in}
\addtolength{\evensidemargin}{-0.375in}
\addtolength{\textwidth}{1in}
\addtolength{\topmargin}{-.45in}
\addtolength{\textheight}{1in}

\urlstyle{rm}

\raggedbottom
\raggedright
\setlength{\tabcolsep}{0in}

\titleformat{\section}{
  \vspace{-10pt}\scshape\raggedright\large
}{}{0em}{}[\color{blue}\titlerule \vspace{-6pt}]

\newcommand{\resumeItem}[2]{
  \item\small{
    \textbf{#1}{: #2 \vspace{-2pt}}
  }
}

\newcommand{\resumeItemWithoutTitle}[1]{
  \item\small{
    {\vspace{-2pt}}
  }
}

\newcommand{\resumeSubheading}[4]{
  \vspace{-1pt}\item
    \begin{tabular*}{0.97\textwidth}{l@{\extracolsep{\fill}}r}
      \textbf{#1} & #2 \\
      \textit{#3} & \textit{#4} \\
    \end{tabular*}\vspace{-5pt}
}

\usepackage{array}
\newenvironment{awlist}{%
    \begin{tabular}{>{\bfseries}l @{\hspace{1em}} p{13cm}}
}{%
    \end{tabular}
}

\newcommand{\resumeSubItem}[2]{\resumeItem{#1}{#2}\vspace{-3pt}}

\renewcommand{\labelitemii}{$\circ$}

\newcommand{\resumeSubHeadingListStart}{\begin{itemize}[leftmargin=*]}
\newcommand{\resumeSubHeadingListEnd}{\end{itemize}}
\newcommand{\resumeItemListStart}{\begin{itemize}}
\newcommand{\resumeItemListEnd}{\end{itemize}\vspace{-5pt}}

\usepackage{fontawesome5}
\usepackage{hyperref}
\hypersetup{colorlinks=true, urlcolor=blue}

\usepackage{charter}

%-----------------------------
%%%%%%  CV STARTS HERE  %%%%%%

\begin{document}

%----------HEADING-----------------

\begin{center}
    \textbf{{\LARGE Harsh Bordekar}} \\
    \textbf{AI Scientist - Physics Models} \\
    \textbf{Physics-Informed ML for Simulation, Surrogate Modelling \& Uncertainty Quantification} \\
    \textbf{PhD Candidate, DLR \& TU Braunschweig | Aerospace Composites \& Structural Simulation}

    \vspace{0.5em}

    \small
    \begin{tabular}{*{3}{l@{\hspace{2.5em}}}}
        \faIcon[regular]{envelope} \href{mailto:harshbordekar@gmail.com}{harshbordekar@gmail.com} &
        \faIcon{linkedin} \href{https://www.linkedin.com/in/harsh-s-bordekar/}{Harsh's LinkedIn} &
        \faIcon{github} \href{https://github.com/harshbordekar-stack}{github.com/harshbordekar-stack} \\
        \faIcon{mobile-alt} +49 (0) 17673888310 &
        \faIcon{map-marker-alt} Braunschweig, Germany &
        \faIcon{id-card} Permanent Resident - Germany
    \end{tabular}
\end{center}

\vspace{5pt}

\section{\color{blue}Profile}
Simulation and machine-learning research engineer with 7+ years developing physics-based models for aerospace structures, including 3+ years applying ML and surrogate modelling directly to simulation problems. PhD work builds a probabilistic, physics-based multiscale surrogate framework for CFRP composites, coupling microstructure-informed mesoscale damage models with stochastic spatial defect distributions to propagate failure from ply to structural scale and replacing computationally expensive full-scale FEM with fast, physics-consistent predictions. The framework is evaluated against a certification-aligned failure criterion derived as part of the same work. Runs large-scale simulation campaigns on HPC/SLURM using domain-specific solvers (ABAQUS, Nastran/Patran) to generate the high-fidelity data these models depend on, and applies Monte Carlo uncertainty quantification to characterise prediction confidence ahead of design decisions. This is direct, hands-on experience with the UQ/out-of-distribution estimation this role treats as a first-class concern. Also built a physics-based CNN surrogate for on-demand computational homogenization, exposed through a tool-calling interface, replacing full-scale FEM homogenization runs with fast inference. Comfortable across the applied-ML stack (PyTorch, TensorFlow, scikit-learn) and in Python/Linux/HPC environments day to day, with a track record of verifying and technically reviewing externally produced simulation models against defined acceptance criteria, and of explaining simulation and ML concepts to both engineering and non-technical audiences.

\vspace{0.2cm}
\textbf{Core Competencies:} ML for Simulation \& Surrogate Modelling, Physics-Informed Machine Learning, Multiscale / Multi-Fidelity Damage Modelling, Uncertainty Quantification \& Out-of-Distribution Estimation, High-Fidelity Simulation Data Generation, FE Solvers (ABAQUS, Nastran/Patran), HPC-Scale Simulation Campaigns (SLURM), Model Verification \& Validation, Physics-Based CNN / Surrogate Modelling, Explainable AI, Python

\section{\color{blue}Professional Experience}

\resumeSubheading
    {Structural Analysis / Simulation Engineer}{DLR \& TU Braunschweig}
    {Physics-informed ML for simulation \& surrogate modelling (PhD)}{Feb $2021$ - Present}
    \resumeItemListStart

    \item (PhD) Developed a probabilistic, physics-based multiscale surrogate framework for leakage-onset prediction in CFRP hydrogen storage tanks, coupling microstructure-informed mesoscale damage models with stochastic spatial defect distributions to propagate failure across ply, laminate, and structural scales, replacing full-scale FEM runs with a physics-consistent surrogate and estimating probability of failure directly from the simulation-generated dataset.
    \\
    \textbf{Tools: Python, ABAQUS, FEM, HPC/SLURM, Multiscale Modelling}

    \item (PhD) Derived a certification-aligned leakage failure criterion mapping combined strain/load states to leakage onset, bridging progressive-damage models and system-level tank simulations and providing the physics-based evaluation standard against which the surrogate framework's coverage, accuracy, and failure modes are assessed.
    \\
    \textbf{Tools: ABAQUS, Python, FEM post-processing using C++}

    \item Performed Monte Carlo uncertainty quantification over the multiscale simulation outputs to generate leakage-probability-vs-strain relationships, characterising prediction confidence and enabling risk-informed design decisions for next-generation CFRP pressure vessels. This is direct UQ/out-of-distribution experience applied to physics-simulation outputs.
    \\
    \textbf{Tools: Python (SciPy, NumPy), Monte Carlo methods, ABAQUS}

    \item Ran large-scale thermal and structural simulation campaigns on HPC infrastructure to evaluate thermal-protection effectiveness for the Callisto reusable-rocket programme's metal upper housing bay, using domain-specific solvers to generate high-fidelity simulation data under production-scale compute constraints.
    \\
    \textbf{Tools: ABAQUS, SLURM (HPC), Python, Nastran, Patran}

    \item Built a physics-based CNN surrogate for microstructure-informed computational homogenization, wrapped in an agentic, tool-calling interface that accepts natural-language queries and microstructure images, retrieves context via a hierarchical RAG pipeline over \textbf{50} materials documents, and returns homogenized/effective material properties, replacing full-scale FEM homogenization runs (hours) with on-demand inference (seconds). Retrieval/prediction faithfulness validated at \textbf{93}\% against ground-truth homogenization results on a held-out set of more than \textbf{100} cases.
    \\
    \textbf{Tools: Python, Physics-Based CNN, Hierarchical RAG, LLM Agent Framework, Image Processing}

    \item Verified and technically reviewed externally and student-produced simulation models, meshes, and results against defined methods and acceptance criteria, diagnosing failure modes arising from data gaps or modelling limitations and issuing documented sign-off-style findings.
    \\
    \textbf{Tools: ABAQUS, Nastran, Python}

    \item Instructed 60+ Master's students in FEM theory and ABAQUS-based CFRP modelling, covering Continuum Damage Mechanics (CDM), Progressive Damage Analysis (PDA), and Cohesive Zone Modelling (CZM). This demonstrated the ability to communicate simulation and ML concepts to both technical and non-technical audiences.
    \\
    \textbf{Tools: ABAQUS, Automation using Python}

    \resumeItemListEnd

    \vspace{10pt}

\resumeSubheading
    {Research Intern, Metallic Fatigue \& Damage Tolerance}{DLR-Braunschweig}
    {Explainable AI \& physics-based fatigue/defect-tolerance assessment}{Nov $2019$ - Jan $2021$}
    \resumeItemListStart

    \item Developed a CT-based, explainable-AI defect-detection framework achieving \textbf{91}\% classification accuracy (\textbf{89}\%  precision / \textbf{91}\%  recall) on \textbf{1200} CT scans of additively manufactured titanium parts, cutting manual review time to \textbf{$<$0.5 sec per image}, with SHAP, Grad-CAM and LIME outputs used to support engineering sign-off. \textit{(Published: Journal of Intelligent Manufacturing, 2025.)}
    \\
    \textbf{Tools: Python (scikit-learn, TensorFlow), SHAP, Grad-CAM, LIME, CT image processing}

    \item Assessed fatigue and defect tolerance of aerospace-grade titanium components using small-defect fatigue models (Murakami, El-Haddad), coupling defect measurements with numerical simulation to translate physics-based fracture-mechanics criteria into allowable-defect acceptance limits.
    \\
    \textbf{Tools: Small-defect fatigue models (Murakami, El-Haddad), Python (NumPy/SciPy), FEM (ABAQUS)}

    \resumeItemListEnd

    \vspace{10pt}

\resumeSubheading
    {Student Research Assistant}{Leibniz University Hannover}
    {Institute of Static and Dynamic}{Sep $2018$ - Oct $2019$}
    \resumeItemListStart

    \item Developed an ML framework for automated detection and quantification of fiber undulations in unidirectional CFRP laminates from image data, generating spatially resolved defect metrics used directly as high-fidelity inputs for probabilistic structural simulation models assessing manufacturing-induced variability.
    \\
    \textbf{Tools: MATLAB, Python (scikit-learn, TensorFlow), image processing, NumPy/SciPy}

    \resumeItemListEnd

%-------------SKILLS SUMMARY------------------
\renewcommand{\arraystretch}{1.3}
\section{\color{blue}Skills Summary}
\begin{tabular}{|p{8cm}p{10.5cm}|}

\hline
\textbf{ML for Simulation \& Surrogate Modelling} &
Physics-Informed Machine Learning \(\vert\) Multiscale / Multi-Fidelity Damage Modelling \(\vert\) Physics-Based CNN Surrogates \(\vert\) Computational Homogenization \\
\hline

\textbf{Uncertainty \& Evaluation} &
Uncertainty Quantification (Monte Carlo) \(\vert\) Verification \& Validation \(\vert\) Physics-Based Evaluation Against Certification Criteria \(\vert\) Failure-Mode Diagnosis \\
\hline

\textbf{Simulation Solvers \& HPC} &
ABAQUS \(\vert\) Nastran/Patran \(\vert\) Linear \& Non-linear FEM \(\vert\) HPC / SLURM (Large-Scale Simulation Campaigns) \(\vert\) Analytical Stress Methods (Niu, Bruhn) \\
\hline

\textbf{ML \& AI} &
PyTorch \(\vert\) TensorFlow \(\vert\) Scikit-learn \(\vert\) Explainable AI (SHAP, Grad-CAM, LIME) \(\vert\) Hugging Face Transformers (Fine-tuning) \\
\hline

\textbf{LLMs \& Agentic Systems} &
Retrieval-Augmented Generation (RAG) \(\vert\) LangChain \(\vert\) LangGraph \(\vert\) Model Context Protocol (MCP) \(\vert\) Tool-Calling \\
\hline

\textbf{Programming Languages} &
Python \(\vert\) C++ \(\vert\) FORTRAN \(\vert\) MATLAB \\
\hline

\textbf{Data Analysis \& Scientific Computing} &
NumPy \(\vert\) Pandas \(\vert\) Matplotlib \(\vert\) Seaborn \(\vert\) Plotly \(\vert\) SALib \(\vert\) SciPy \\
\hline

\textbf{APIs \& Backend Development} &
RESTful APIs \(\vert\) FastAPI \(\vert\) JSON \\
\hline

\textbf{Version Control \& Reproducibility} &
Git \(\vert\) Docker \(\vert\) Linux \\
\hline

\textbf{Application Development} &
VS Code \(\vert\) JupyterLab \(\vert\) Google Colab \(\vert\) Streamlit \(\vert\) Gradio \\
\hline

\textbf{Languages} &
English (C1) \(\vert\) German (B1) \\
\hline
\end{tabular}

\vspace{0.2cm}
\vspace{5pt}

\newpage
%-----------Awards-----------------

%-----------EDUCATION-----------------
\section{\color{blue}Awards and Education}

\resumeSubheading
  {\textcolor{green!50!black}{\faAward} \textbf{2nd Prize: Best Paper Award}}
  {DLR (German Aerospace Center)}{}{2026}
\resumeSubheading
  {\textcolor{green!50!black}{\faAward} \textbf{3rd Prize: Best Paper Award}}
  {DLR (German Aerospace Center)}{}{2024}
\vspace{4pt}

\resumeSubheading
  {\textcolor{green!50!black}{PhD Candidate,} Computational Methods in Engineering}
  {DLR \& TU Braunschweig, Germany}{}{Feb 2021 -- 2026 (expected)}
  \vspace{2pt}
\textbf{Dissertation:} Physics-Informed, Multiscale Machine Learning for CFRP Composite Structures
\vspace{4pt}

\resumeSubheading
  {\textcolor{green!50!black}{M.Sc.} Computational Methods in Engineering}
  {Leibniz University Hannover, Germany}{}{April 2018 -- Dec 2020}
  \vspace{2pt}
\textbf{Focus:} Numerical Methods, Scientific Machine Learning, Deep Learning
\vspace{4pt}

\section{\color{blue}Selected Publications}

\begin{enumerate}

  \item \textbf{Bordekar, H.}, et al.
  ``Microstructure-informed Probabilistic Multiscale Modelling of CFRP via Machine-Learning-Derived Stochastic Volume Elements.''
  \textit{Manuscript in preparation}, 2026.

  \item \textbf{Bordekar, H.}, et al.
  ``MicroStructAgent: A Physics-Informed Agentic AI Framework for Computational Homogenization and Quality Assurance of Fiber Composite Laminates.''
  \textit{Manuscript in preparation}, 2026.

  \item \textbf{Bordekar, H.}, et al.
  ``Microstructure Characterization of Fiber Composite Laminate Using eXplainable Artificial Intelligence.''
  \textit{Journal of Composite Materials}, vol. 60, no. 1, 2026, pp. 3--25.
  \href{https://journals.sagepub.com/doi/pdf/10.1177/00219983251345559}{\textcolor{blue}{DOI}}

  \item \textbf{Bordekar, H.}, et al.
  ``eXplainable Artificial Intelligence for Automatic Defect Detection in Additively Manufactured Parts Using CT Scan Analysis.''
  \textit{Journal of Intelligent Manufacturing}, vol. 36, no. 2, 2025, pp. 957--974.
  \href{https://link.springer.com/article/10.1007/s10845-023-02272-4}{\textcolor{blue}{DOI}}

\end{enumerate}
\vspace{2pt}

\section{\color{blue}References}

\begin{tabular}{@{} p{6cm} p{10cm}@{}}
\textbf{\href{https://www.dlr.de/de/sy/ueber-uns/abteilungen-dlr-sy/dlr-sy-funktionsleichtbau}{Prof. Dr. Christian Huehne}} & Abteilungsleiter Funktionsleichtbau. Deutsches Zentrum für Luft- und Raumfahrt. Institut für Systemleichtbau \\
\textbf{\href{https://www.tu-braunschweig.de/ima/team/head-of-institute/prof-dr-ing-oliver-voelkerink}{Prof. Dr. Oliver Voelkerink}} & Head of Institute., IMA, TU-Braunschweig\\
\textbf{\href{https://www.linkedin.com/in/nicola-cersullo-41772b131 }{Dr. Nicola Cersullo}} & R\&T Project Manager @ Airbus Hamburg \\
\end{tabular}

\end{document}
